(1) 代入 \(\xi=x-vt-x_0\) 后得到 \(-vU'+6UU'+U'''=0\)。第一次积分并使用衰减边界条件，得
\[
U''-vU+3U^2=0.
\]
乘以 \(U'\) 再积分，得到 \((U')^2=vU^2-2U^3=U^2(v-2U)\)。正的非平凡局域波要求 \(v>0\)，在波峰 \(U'=0\) 时 \(U_{\max}=v/2\)。直接代入或分离变量积分可得唯一的平移族
\[
\boxed{U(\xi)=\frac v2\operatorname{sech}^2\!\left(\frac{\sqrt v}{2}\xi\right)}.
\]
令 \(k=\sqrt v/2\)，则宽度量级 \(1/k=2/\sqrt v\)：振幅越高，传播越快而宽度越窄。将该表达式代回 \(U''-vU+3U^2=0\) 可逐项核验，且两端确实指数衰减。

(2) 方程写成守恒形式 \(u_t+\partial_x(3u^2+u_{xx})=0\)，故 \(\dot M=0\)。乘以 \(2u\) 后有
\[
\partial_t(u^2)+\partial_x(4u^3+2uu_{xx}-u_x^2)=0,
\]
因而 \(\dot I=0\)。使用 \(\int_{-\infty}^{\infty}\operatorname{sech}^2(k\xi)d\xi=2/k\) 与 \(\int\operatorname{sech}^4(k\xi)d\xi=4/(3k)\)，得到
\[
\boxed{M=\frac v2\frac2k=2\sqrt v},\qquad
\boxed{I=\frac{v^2}{4}\frac4{3k}=\frac23v^{3/2}}.
\]
二者均正且随波速单调增大；若 \(v\to0^+\)，孤立波振幅、质量和二次积分均趋零，而宽度趋于无穷大。

(3) 变分给出 \(\delta H/\delta u=3u^2+u_{xx}\)，因此 \(u_t=-\partial_x(\delta H/\delta u)\)。于是
\[
\frac{dH}{dt}=\int\frac{\delta H}{\delta u}\,u_t\,dx
=-\frac12\int\partial_x\!\left(\frac{\delta H}{\delta u}\right)^2dx=0.
\]
对孤立波使用 \(\int\operatorname{sech}^6(k\xi)d\xi=16/(15k)\)，并由 (1) 的第一积分得到
\[
\int U^3dx=\frac4{15}v^{5/2},\qquad
\int(U')^2dx=vI-2\int U^3dx=\frac2{15}v^{5/2}.
\]
故
\[
\boxed{H=\frac4{15}v^{5/2}-\frac12\frac2{15}v^{5/2}=\frac15v^{5/2}}.
\]
梯度项虽为负号，却是保证变分动力学正确色散项 \(u_{xxx}\) 的必要部分；不能把 \(H\) 简单解释为处处非负的机械能密度。

(4) 若 \(u(x,t)\) 是解，则 \(u_{u_0}(x,t)=u(x-6u_0t,t)+u_0\) 仍是解：时间导数额外产生 \(-6u_0u_x\)，恰好抵消非线性项新增的 \(6u_0u_x\)。因而背景 \(u_0\) 上的孤立波相对背景振幅仍为 \(v/2\)，实验室传播速度变为 \(v+6u_0\)。另一个对称是
\[
u_\lambda(x,t)=\lambda^2u(\lambda x,\lambda^3t),\qquad\lambda>0.
\]
各项均乘 \(\lambda^5\)，因此变换保持方程；\(M,I,H\) 分别按 \(\lambda,\lambda^3,\lambda^5\) 缩放，与 (2)(3) 的 \(v\mapsto\lambda^2v\) 一致。在线性小振幅极限，取 \(u\propto e^{i(kx-\omega t)}\)，得
\[
-i\omega+(ik)^3=0,\quad \boxed{\omega=-k^3},\quad
\boxed{v_g=\frac{d\omega}{dk}=-3k^2\le0}.
\]
所以在零背景上小振幅线性波向负方向输运，而非线性孤立波 \(v>0\) 向正方向传播；非线性与色散的平衡使这种传播方向差异成为可能。
